% ============================================================
% Bibby AI — Default Article Template (pdflatex)
% Compile: pdflatex -> biber -> pdflatex -> pdflatex
% ============================================================
\documentclass[11pt]{article}

% --- Core ---
\usepackage[utf8]{inputenc}
\usepackage[T1]{fontenc}
\usepackage[margin=1in]{geometry}
\usepackage{microtype}

% --- Math ---
\usepackage{amsmath, amssymb, amsthm, mathtools}

% --- Floats ---
\usepackage{graphicx}
\usepackage{booktabs}
\usepackage{subcaption}
\usepackage{siunitx}

% --- Bibliography ---
\usepackage[backend=biber, style=numeric-comp, sorting=none]{biblatex}
\addbibresource{references.bib}

% --- Links & refs (order matters: hyperref then cleveref, both last) ---
\usepackage[colorlinks=true, linkcolor=blue, citecolor=blue, urlcolor=blue]{hyperref}
\usepackage[capitalize]{cleveref}

% --- Theorem environments ---
\newtheorem{theorem}{Theorem}[section]
\newtheorem{lemma}[theorem]{Lemma}
\newtheorem{proposition}[theorem]{Proposition}
\newtheorem{corollary}[theorem]{Corollary}
\theoremstyle{definition}
\newtheorem{definition}[theorem]{Definition}
\theoremstyle{remark}
\newtheorem{remark}[theorem]{Remark}

% --- Semantic macros ---
\newcommand{\R}{\mathbb{R}}
\newcommand{\N}{\mathbb{N}}
\newcommand{\Z}{\mathbb{Z}}
\newcommand{\E}{\mathbb{E}}
\newcommand{\Prob}{\mathbb{P}}
\newcommand{\norm}[1]{\left\lVert #1 \right\rVert}
\newcommand{\abs}[1]{\left\lvert #1 \right\rvert}
\newcommand{\inner}[2]{\left\langle #1, #2 \right\rangle}
\newcommand{\given}{\mid}
\DeclareMathOperator*{\argmin}{arg\,min}
\DeclareMathOperator*{\argmax}{arg\,max}

% --- Metadata ---
\title{Paper Title}
\author{First Author\thanks{Affiliation, email} \and Second Author}
\date{\today}

\begin{document}

\maketitle

\begin{abstract}
One paragraph. Problem, approach, key result, significance.
\end{abstract}

\section{Introduction}\label{sec:intro}
Motivation and contributions.
As shown in \cref{sec:method}, our approach improves on prior work~\cite{examplekey2024}.

\section{Method}\label{sec:method}
Given data $x \in \R^d$, we minimize
\begin{equation}\label{eq:objective}
  \theta^\star = \argmin_{\theta} \; \E\!\left[ \mathcal{L}(x; \theta) \right].
\end{equation}

\section{Experiments}\label{sec:experiments}
See \cref{tab:results} and \cref{eq:objective}.

\begin{table}[t]
  \centering
  \caption{Main results.}
  \label{tab:results}
  \begin{tabular}{lcc}
    \toprule
    Method & Metric A & Metric B \\
    \midrule
    Baseline & 0.0 & 0.0 \\
    Ours     & \textbf{0.0} & \textbf{0.0} \\
    \bottomrule
  \end{tabular}
\end{table}

\section{Conclusion}\label{sec:conclusion}
Summary and future work.

\printbibliography

\end{document}